% A short paper with what a reader gets around by: an outline, links inside
% (equations, a figure, sections, citations) and one to the web. Typeset with
% `pdflatex around.tex` twice; `ILMI_SMOKE=sioyek` reads around.pdf.
\documentclass[11pt]{article}
\usepackage[a5paper,margin=18mm]{geometry}
\usepackage{amsmath}
\usepackage[colorlinks=false,pdfborder={0 0 0},bookmarksnumbered=true]{hyperref}
\title{Fluctuations in the canonical ensemble}
\author{H. Lindqvist}
\date{}
\begin{document}
\maketitle
\tableofcontents

\section{The partition function}\label{sec:z}
A system held at temperature $T$ by a reservoir occupies a microstate $s$ of
energy $E_s$ with probability
\begin{equation}\label{eq:p}
  P_s = \frac{e^{-\beta E_s}}{Z}, \qquad Z = \sum_s e^{-\beta E_s},
\end{equation}
where $\beta = 1/kT$. Everything below follows from $Z$ by differentiation;
the entropy of the reservoir has already been spent in deriving
\eqref{eq:p}, as Gibbs first argued~\cite{gibbs}. The mean energy is
\begin{equation}\label{eq:u}
  U = -\frac{\partial \ln Z}{\partial \beta}.
\end{equation}
The sketch in Fig.~\ref{fig:well} shows the two levels used in
Section~\ref{sec:two}; the heat capacity of that system is derived there
and the general result in Section~\ref{sec:fluct}.

\newpage
\section{Energy fluctuations}\label{sec:fluct}
Differentiating \eqref{eq:u} once more gives the variance of the energy,
\begin{equation}\label{eq:var}
  \sigma_E^2 = \frac{\partial^2 \ln Z}{\partial \beta^2} = kT^2 C_V ,
\end{equation}
so a heat capacity is a measure of fluctuation. The entropy is extensive,
and so is $C_V$; the relative width $\sigma_E / U$ therefore falls as
$N^{-1/2}$. Einstein used exactly this to argue that entropy is
statistical~\cite{einstein}. A fuller account is in Pathria's
text~\cite{pathria}, and a readable summary is kept at
\url{https://en.wikipedia.org/wiki/Canonical_ensemble}.

\begin{figure}[b]
\centering
\setlength{\unitlength}{1mm}
\begin{picture}(60,24)
\put(5,4){\line(1,0){50}}
\put(5,18){\line(1,0){50}}
\put(30,4){\vector(0,1){14}}
\put(57,3){$0$}
\put(57,17){$\epsilon$}
\end{picture}
\caption{The two-level system: a ground state and one excited state at
energy $\epsilon$.}\label{fig:well}
\end{figure}

\newpage
\section{A two-level system}\label{sec:two}
For the system of Fig.~\ref{fig:well} the sum has two terms,
$Z = 1 + e^{-\beta\epsilon}$, and \eqref{eq:u} gives
\begin{equation}\label{eq:two}
  U = \frac{\epsilon}{e^{\beta\epsilon} + 1}.
\end{equation}
Its heat capacity, from \eqref{eq:var}, rises to a peak near
$kT \approx 0.42\,\epsilon$ and falls again: the Schottky anomaly. The
entropy climbs from zero to $k \ln 2$ across the same range.

\newpage
\section{Many systems}
Independent systems multiply their partition functions, so their
entropy and their energy add, and \eqref{eq:var} adds too. This is why
the fluctuations of Section~\ref{sec:fluct} are invisible in bulk matter.

\begin{thebibliography}{9}
\bibitem{gibbs} J.~W. Gibbs, \emph{Elementary Principles in Statistical
Mechanics}. Scribner, 1902.
\bibitem{einstein} A. Einstein, ``Zur allgemeinen molekularen Theorie der
W\"arme.'' \emph{Annalen der Physik} 14, 1904.
\bibitem{pathria} R.~K. Pathria, \emph{Statistical Mechanics}, 2nd ed.
Butterworth-Heinemann, 1996.
\end{thebibliography}
\end{document}
